
\documentclass[11pt,a4paper]{article}
\usepackage[utf8]{inputenc}
\usepackage[T1]{fontenc}
\usepackage[french]{babel}
\usepackage{lmodern}
\usepackage{microtype}
\usepackage[a4paper,margin=1.7cm,headheight=15pt]{geometry}
\usepackage{amsmath,amssymb,mathtools}
\usepackage{array,tabularx,multicol}
\usepackage{enumitem}
\usepackage[most]{tcolorbox}
\usepackage{xcolor}
\usepackage{tikz}
\usepackage{pgfplots}
\pgfplotsset{compat=1.18}
\usetikzlibrary{arrows.meta,calc,patterns}
\usepackage{fancyhdr}
\usepackage{fancyvrb}
\usepackage{listings}
\usepackage{hyperref}
\usepackage{pdfpages}

\definecolor{bleufonce}{RGB}{28,55,120}
\definecolor{bleuclair}{RGB}{238,243,255}
\definecolor{grisclair}{RGB}{246,246,246}
\definecolor{tableline}{RGB}{45,45,45}
\definecolor{softgray}{RGB}{246,247,249}
\definecolor{deepblue}{RGB}{25,62,115}
\definecolor{softblue}{RGB}{235,241,250}
\definecolor{accent}{RGB}{20,82,130}

\hypersetup{colorlinks=true,linkcolor=bleufonce,urlcolor=bleufonce,pdfauthor={},pdftitle={Parcours de révision -- Trigonométrie et LGN -- Énoncés}}
\setlength{\parindent}{0pt}
\setlength{\parskip}{0.28em}
\renewcommand{\arraystretch}{1.12}
\setlist[itemize]{leftmargin=1.25em,itemsep=0.15em,topsep=0.2em}
\setlist[enumerate,1]{label=\textbf{\arabic*.}, leftmargin=1.05cm, itemsep=0.35em, topsep=0.25em}
\setlist[enumerate,2]{label=\textbf{\alph*.}, leftmargin=0.85cm, itemsep=0.25em, topsep=0.2em}

\pagestyle{fancy}
\fancyhf{}
\lhead{}
\rhead{}
\cfoot{\thepage}

\newcommand{\e}{\mathrm e}
\newcommand{\dd}{\,\mathrm d}
\newcommand{\vect}[1]{\overrightarrow{#1}}
\newcommand{\origine}[1]{ {\scriptsize\textcolor{bleufonce}{(site #1)}}}
\newcommand{\code}[1]{\texttt{\detokenize{#1}}}
\newcommand{\Bin}{\mathcal B}
\newcounter{question}
\newcounter{reponse}

\newenvironment{blocquestions}[1]{%
  \begin{tcolorbox}[
    enhanced,
    breakable,
    colback=bleuclair,
    colframe=bleufonce,
    coltitle=white,
    colbacktitle=bleufonce,
    title={\bfseries #1},
    fonttitle=\bfseries,
    boxrule=0.6pt,
    arc=2mm,
    left=2mm,
    right=2mm,
    top=2mm,
    bottom=1.5mm,
    before skip=3mm,
    after skip=3mm, fontupper=\large
  ]
  \large
  \begin{enumerate}[leftmargin=*,label=\textbf{\arabic*.},itemsep=0.45em,topsep=0.4em]
  \setcounter{enumi}{\value{question}}
}{%
  \setcounter{question}{\value{enumi}}
  \end{enumerate}
  \end{tcolorbox}
}

\newenvironment{blocreponses}[1]{%
  \begin{tcolorbox}[
    enhanced,
    breakable,
    colback=grisclair,
    colframe=bleufonce,
    coltitle=white,
    colbacktitle=bleufonce,
    title={\bfseries #1},
    fonttitle=\bfseries,
    boxrule=0.6pt,
    arc=2mm,
    left=2mm,
    right=2mm,
    top=2mm,
    bottom=1.5mm,
    before skip=3mm,
    after skip=3mm, fontupper=\large
  ]
  \large
  \begin{enumerate}[leftmargin=*,label=\textbf{\arabic*.},itemsep=0.45em,topsep=0.5em]
  \setcounter{enumi}{\value{reponse}}
}{%
  \setcounter{reponse}{\value{enumi}}
  \end{enumerate}
  \end{tcolorbox}
}

\newtcolorbox{correctionbox}[1]{
  enhanced,
  breakable,
  colback=grisclair,
  colframe=black!55,
  boxrule=0.55pt,
  arc=2mm,
  left=2mm,right=2mm,top=1.5mm,bottom=1.5mm,
  title=\textbf{#1},
  coltitle=black,
  colbacktitle=black!12,
  before skip=2mm,
  after skip=2mm
}

\newcommand{\titreSujet}[2]{%
  \subsection*{#1}
  \addcontentsline{toc}{subsection}{#1}
  \textit{Référence : #2}\par\medskip\hrule\medskip
}

\lstset{language=Python,basicstyle=\ttfamily\small,numbers=left,numberstyle=\tiny,frame=single,showstringspaces=false,columns=fullflexible,keepspaces=true}
\sloppy

\begin{document}
\begin{center}
{\LARGE\bfseries Parcours de révision -- Trigonométrie et loi des grands nombres}\\[1mm]

\end{center}
\vspace{1mm}
\tableofcontents
\clearpage

\phantomsection
\addcontentsline{toc}{section}{Partie 1 -- Récapitulatif de cours}
\includepdf[pages=-,pagecommand={},fitpaper=true]{parcours_revision_trigo_lgn_carte_mentale_illustre_v3.pdf}

\clearpage
\section*{Partie 2 -- Questions flashs}
\addcontentsline{toc}{section}{Partie 2 -- Questions flashs}
\setcounter{question}{0}
\begin{blocquestions}{A. Trigonométrie : cercle et propriétés}
  \item \origine{479} Qu'est-ce que le cercle trigonométrique ?
  \item \origine{480} Que signifie que les fonctions cosinus et sinus sont $2\pi$-périodiques ?
  \item \origine{481} Quelle propriété de parité vérifie la fonction cosinus ?
  \item \origine{482} Quelle propriété de parité vérifie la fonction sinus ?
  \item \origine{528} Quel est le point du cercle trigonométrique associé à un réel $x$ ?
  \item \origine{529} Quelle longueur d'arc correspond à un tour complet sur le cercle trigonométrique ?
  \item \origine{530} Que vaut $\cos^2 x+\sin^2 x$ ?
\end{blocquestions}
\begin{blocquestions}{B. Trigonométrie : valeurs remarquables}
  \item \origine{531} Que valent $\cos(0)$ et $\sin(0)$ ?
  \item \origine{533} Que valent $\cos\left(\dfrac{\pi}{2}\right)$ et $\sin\left(\dfrac{\pi}{2}\right)$ ?
  \item \origine{484} Que valent $\cos\left(\dfrac{\pi}{4}\right)$ et $\sin\left(\dfrac{\pi}{4}\right)$ ?
  \item \origine{691} Donner $\cos\left(\dfrac{\pi}{6}\right)$ et $\sin\left(\dfrac{\pi}{6}\right)$.
\end{blocquestions}
\begin{blocquestions}{C. Liens avec l'analyse}
  \item \origine{158} Donner une primitive de $\cos(x)$.
  \item \origine{159} Donner une primitive de $\sin(x)$.
  \item \origine{247} Quelle est une primitive de $-\sin x$ ?
  \item \origine{798} Déterminer $\displaystyle\lim_{x\to 0}\dfrac{1-\cos x}{x^2}$.
\end{blocquestions}
\begin{blocquestions}{D. Variables aléatoires et loi des grands nombres}
  \item Donner la formule de l'espérance d'une variable aléatoire finie $X$. \origine{489}
  \item Donner la formule de la variance d'une variable aléatoire finie $X$. \origine{490}
  \item Comment définit-on l'écart-type d'une variable aléatoire $X$ ? \origine{491}
  \item Quelle est l'espérance d'une loi binomiale $\mathcal B(n,p)$ ? \origine{63}
  \item Quelle est la variance d'une loi binomiale $\mathcal B(n,p)$ ? \origine{65}
  \item Dans un contexte concret, que représente l'espérance ? \origine{64}
  \item Que représente une moyenne empirique sur un échantillon ?
  \item Que dit, de manière simple, la loi faible des grands nombres ?
  \item Quand $n$ augmente, que devient généralement le risque d'un écart important entre une moyenne empirique et l'espérance ?
  \item Si $X_1,\ldots,X_n$ sont indépendantes et de même espérance $\mu$, que vaut l'espérance de $M_n=\dfrac{X_1+\cdots+X_n}{n}$ ?
\end{blocquestions}

\clearpage
\section*{Partie 3 -- Sujets de bac}
\addcontentsline{toc}{section}{Partie 3 -- Sujets de bac}
\titreSujet{Sujet 1 -- Loi des grands nombres et moyenne d'échantillons}{Métropole, 17 juin 2025, sujet 1, exercice 1}

On compte quatre groupes sanguins dans l'espèce humaine : A, B, AB et O. Chaque groupe sanguin peut présenter un facteur rhésus. Lorsqu'il est présent, on dit que le rhésus est positif, sinon on dit qu'il est négatif.

Au sein de la population française, on sait que :
\begin{itemize}
  \item 45 \% des individus appartiennent au groupe A, et parmi eux 85 \% sont de rhésus positif ;
  \item 10 \% des individus appartiennent au groupe B, et parmi eux 84 \% sont de rhésus positif ;
  \item 3 \% des individus appartiennent au groupe AB, et parmi eux 82 \% sont de rhésus positif.
\end{itemize}

On choisit au hasard une personne dans la population française.

On désigne par :
\begin{itemize}
  \item $A$ l'événement : « la personne choisie est de groupe sanguin A » ;
  \item $B$ l'événement : « la personne choisie est de groupe sanguin B » ;
  \item $AB$ l'événement : « la personne choisie est de groupe sanguin AB » ;
  \item $O$ l'événement : « la personne choisie est de groupe sanguin O » ;
  \item $R$ l'événement : « la personne choisie a un facteur rhésus positif ».
\end{itemize}

Pour un événement quelconque $E$, on note $\overline E$ l'événement contraire de $E$ et $p(E)$ la probabilité de $E$.

\begin{center}
\begin{tikzpicture}[grow=right, level distance=2.8cm, sibling distance=18mm,
  every node/.style={font=\small}, edge from parent/.style={draw},
  edge from parent path={(\tikzparentnode.east) -- (\tikzchildnode.west)}]
\node {}
  child {node {$O$}
    child {node {$\overline R$}}
    child {node {$R$}}
    edge from parent node[below] {$\dots$}}
  child {node {$AB$}
    child {node {$\overline R$} edge from parent node[below] {$\dots$}}
    child {node {$R$} edge from parent node[above] {$\dots$}}
    edge from parent node[below] {$\dots$}}
  child {node {$B$}
    child {node {$\overline R$} edge from parent node[below] {$\dots$}}
    child {node {$R$} edge from parent node[above] {$\dots$}}
    edge from parent node[above] {$\dots$}}
  child {node {$A$}
    child {node {$\overline R$} edge from parent node[below] {$\dots$}}
    child {node {$R$} edge from parent node[above] {$\dots$}}
    edge from parent node[above] {$\dots$}};
\end{tikzpicture}
\end{center}

\begin{enumerate}
  \item Recopier l'arbre ci-dessus en complétant les dix pointillés.
\item Montrer que $p(B\cap R)=0{,}084$. Interpréter ce résultat dans le contexte de l'exercice.
\item On précise que $p(R)=0{,}8397$. Montrer que $p_O(R)=0{,}83$.
\item On dit qu'un individu est « donneur universel » lorsque son sang peut être transfusé à toute personne sans risque d'incompatibilité. Le groupe O de rhésus négatif est le seul vérifiant cette caractéristique. Montrer que la probabilité qu'un individu choisi au hasard dans la population française soit donneur universel est de $0{,}0714$.
\item Lors d'une collecte de sang, on choisit un échantillon de 100 personnes dans la population d'une ville française. Cette population est suffisamment grande pour assimiler ce choix à un tirage avec remise. On note $X$ la variable aléatoire qui à chaque échantillon de 100 personnes associe le nombre de donneurs universels dans cet échantillon.
  \begin{enumerate}
    \item Justifier que $X$ suit une loi binomiale dont on précisera les paramètres.
    \item Déterminer à $10^{-3}$ près la probabilité qu'il y ait au plus 7 donneurs universels dans cet échantillon.
    \item Montrer que l'espérance $E(X)$ de la variable aléatoire $X$ est égale à $7{,}14$ et que sa variance $V(X)$ est égale à $6{,}63$ à $10^{-2}$ près.
  \end{enumerate}
\item Lors de la semaine nationale du don du sang, une collecte de sang est organisée dans $N$ villes françaises choisies au hasard numérotées $1,2,3,\ldots,N$, où $N$ est un entier naturel non nul.

  On considère la variable aléatoire $X_1$ qui à chaque échantillon de 100 personnes de la ville 1 associe le nombre de donneurs universels dans cet échantillon. On définit de la même manière les variables aléatoires $X_2$ pour la ville 2, ..., $X_N$ pour la ville $N$.

  On suppose que ces variables aléatoires sont indépendantes et qu'elles admettent la même espérance égale à $7{,}14$ et la même variance égale à $6{,}63$.

  On considère la variable aléatoire
  \[
  M_N=\dfrac{X_1+X_2+\cdots+X_N}{N}.
  \]
  \begin{enumerate}
    \item Que représente la variable aléatoire $M_N$ dans le contexte de l'exercice ?
    \item Calculer l'espérance $E(M_N)$.
    \item On désigne par $V(M_N)$ la variance de la variable aléatoire $M_N$. Montrer que
    \[
    V(M_N)=\dfrac{6{,}63}{N}.
    \]
    \item Déterminer la plus petite valeur de $N$ pour laquelle l'inégalité de Bienaymé-Tchebychev permet d'affirmer que
    \[
    P(7<M_N<7{,}28)>0{,}95.
    \]
  \end{enumerate}
\end{enumerate}
\clearpage
\titreSujet{Sujet 2 -- Trigonométrie, convexité et intégration}{Amérique du Nord, 22 mai 2025, sujet 2, exercice 4}

On désigne par $f$ la fonction définie sur l'intervalle $[0;\pi]$ par
\[
f(x)=e^x\sin(x).
\]
On note $\mathcal C_f$ la courbe représentative de $f$ dans un repère.

\textbf{Partie A}

\begin{enumerate}
  \item
  \begin{enumerate}
    \item Démontrer que, pour tout réel $x$ de l'intervalle $[0;\pi]$,
    \[
f'(x)=e^x[\sin(x)+\cos(x)].
    \]
    \item Justifier que la fonction $f$ est strictement croissante sur l'intervalle $\left[0;\dfrac{\pi}{2}\right]$.
  \end{enumerate}
  \item
  \begin{enumerate}
    \item Déterminer une équation de la tangente $T$ à la courbe $\mathcal C_f$ au point d'abscisse $0$.
    \item Démontrer que la fonction $f$ est convexe sur l'intervalle $\left[0;\dfrac{\pi}{2}\right]$.
    \item En déduire que, pour tout réel $x$ de l'intervalle $\left[0;\dfrac{\pi}{2}\right]$,
    \[
e^x\sin(x)\geqslant x.
    \]
  \end{enumerate}
  \item Justifier que le point d'abscisse $\dfrac{\pi}{2}$ de la courbe représentative de la fonction $f$ est un point d'inflexion.
\end{enumerate}

\textbf{Partie B}

On note
\[
I=\int_0^{\frac{\pi}{2}}e^x\sin(x)\,dx\qquad\text{et}\qquad J=\int_0^{\frac{\pi}{2}}e^x\cos(x)\,dx.
\]

\begin{enumerate}
  \item En intégrant par parties l'intégrale $I$ de deux manières différentes, établir les deux relations suivantes :
  \[
  I=1+J\qquad\text{et}\qquad I=e^{\frac{\pi}{2}}-J.
  \]
  \item En déduire que
  \[
  I=\dfrac{1+e^{\frac{\pi}{2}}}{2}.
  \]
  \item On note $g$ la fonction définie sur $\mathbb R$ par $g(x)=x$.

  Les courbes représentatives des fonctions $f$ et $g$ sont tracées dans le repère orthogonal ci-dessous sur l'intervalle $[0;\pi]$.
  Calculer la valeur exacte de l'aire du domaine hachuré situé entre les courbes $\mathcal C_f$ et $\mathcal C_g$ et les droites d'équation $x=0$ et $x=\dfrac{\pi}{2}$.
\end{enumerate}

\begin{center}
\begin{tikzpicture}
\begin{axis}[
width=11cm,height=6cm,
axis lines=middle,
xmin=0,xmax=3.25,
ymin=0,ymax=9,
xtick={0,1,2,3},
ytick={0,2,4,6,8},
grid=both,
samples=200,
]
\addplot[domain=0:3.14159, thick] {exp(x)*sin(deg(x))} node[pos=0.80, above] {$\mathcal C_f$};
\addplot[domain=0:3.14159, thick] {x} node[pos=0.90, below] {$\mathcal C_g$};
\addplot[domain=0:1.5708, fill=gray!25, draw=none] {exp(x)*sin(deg(x))} \closedcycle;
\addplot[domain=0:1.5708, fill=white, draw=none] {x} \closedcycle;
\end{axis}
\end{tikzpicture}
\end{center}

\newpage
\end{document}
